\section*{CHAPTER 1: INTRODUCTION}
\addcontentsline{toc}{section}{\numberline{}CHAPTER 1: INTRODUCTION}
\setcounter{section}{1}
\setcounter{subsection}{0}
\setcounter{figure}{0}
\setcounter{table}{0}

\subsection{Problem}

Issue trackers are the primary interface between a software project and its
users. Before any report can be check and assigned priorities, someone must decide what kind of report it is, whether it is a defect, a request for new functionality, or a question about existing behaviour. That determines who sees the issue and how urgently, and it is a repetitive and high-volume work that maintainers mostly perform manually.

Automating it is a classification problem over short, noisy and semi-structured text. Issue reports are Markdown documents in which prose is interleaved with stack traces, code blocks, logs, screenshots and template boilerplate. They vary in length, from a single sentence to over twenty-one
thousand words. And the class boundaries are blurry in which a user who is
unsure whether behaviour is broken and writes something that looks like both a bug report and a question.

\subsection{Task definition}

% Competition page: https://nlbse2024.github.io/tools/
% Official repo (rules, data, baseline): https://github.com/nlbse2024/issue-report-classification
% Overview paper: Kallis et al., "The NLBSE'24 Tool Competition", NLBSE'24,
%   https://doi.org/10.1145/3643787.3648038
We address the \textbf{NLBSE'24 tool competition on issue report classification}~\cite{nlbse2024}.
Compared with the 2023 edition, the organisers deliberately shifted the task towards
few-shot learning on a small multi-repository dataset.
% Source (tools page, "Updates"): "focus on few-shot learning, small multi repo dataset,
%   SentenceTransformer baseline" -- https://nlbse2024.github.io/tools/

\noindent The dataset has 3,000 labelled issue reports from five open-source projects:
\texttt{facebook/react}, \texttt{tensorflow/tensorflow}, \texttt{microsoft/vscode}, \texttt{bitcoin/bitcoin} and \texttt{opencv/opencv}.
Each report carries four fields (repository, label, title and body) and exactly one
label from \texttt{\{bug, feature, question\}}.
% Source (README, "Dataset"): fields = Repository, Label, Title, Body; labels = bug, feature, question
%   https://github.com/nlbse2024/issue-report-classification
% Data files: https://github.com/nlbse2024/issue-report-classification/tree/main/data
%   (issues_train.csv, issues_test.csv)


\begin{itemize}
    \item \textbf{The task is multi-class, not multi-label.} The organisers excluded
    every issue carrying more than one label, so each report belongs to exactly one class,
    and each classifier must assign exactly one label per issue.
    % Source (README): "Issues with multiple labels are excluded from the dataset."
    %   and "The classifier should assign one label to an issue."
    
    \item \textbf{Five classifiers are required.} The process trains a separate
    model for each repository and evaluates it only on that repository's test issues.
    Each classifier therefore sees only 300 training examples (100 per class), this is a few-shot problem by construction.
    % Source (README, "Training"): "Participants must train a multi-class classifier for
    %   each of the 5 projects, i.e., we expect 5 classifiers per model submission."
    % 300 per repo = 1,500 training issues / 5 repos; 100 per class follows from balance.

    \item \textbf{All training data must come from the provided training set.}
    Pretrained models are permitted, but they may only be fine-tuned on the given
    training data, and any input or feature must be derived from it. Preprocessing,
    feature engineering, resampling and carving out a validation split are allowed
    on the training set. On the test set, any preprocessing is allowed except
    sampling or rebalancing.
    % Source (README, "Training" and "Evaluation" sections)
\end{itemize}

\noindent The data is split 50/50 into a training set (1,500 issues) and a test set
(1,500 issues). Both are class-balanced, which works out to 100 issues per
(project, class) cell in each split.
% Source (README): "split into a training set (50%) and a test set (50%)";
%   training set "encompassing 1.5 thousand labeled issue reports"; both sets balanced.
% 100 per cell = 1,500 / (5 repos x 3 classes). Not stated verbatim by the organisers,
%   but consistent with the baseline's per-class recalls (multiples of 0.01).

\noindent Per-repository performance is the macro-averaged F1 over the three classes. The cross-repository score is the arithmetic mean of the five per-repository F1 scores, and submissions are ranked on this score alone. Submissions must also report precision, recall and F1 for each class and each repository.
% Source (README, "Evaluation"): "Repository classification performance is measured using
%   the F1 score over all the three classes (averaged). Cross-repository performance is
%   measured as the arithmetic mean of the five repository F1 scores."
% Note: the README says macro and micro give the same result on balanced data. That holds
%   for recall (both equal accuracy) but not for F1. The baseline table uses macro:
%   react macro-F1 = 0.8718 vs accuracy/micro-F1 = 0.8733.

\subsection{Approach}

We evaluate twenty-two models spanning the range of available approaches, all measured through a single evaluation setup so that every leaderboard entry is computed identically:

% \begin{itemize}
%     \item \textbf{Reference floors} - majority class, stratified random, hand-written keyword rules, and naive Bayes implemented from scratch.
%     \item \textbf{Classical machine learning} - TF-IDF features with naive Bayes, logistic regression, linear SVM and random forest, with a preprocessing ablation and hyperparameter search.
%     \item \textbf{Neural networks} - a feed-forward network over TF-IDF features and a 1-D convolutional network over learned embeddings.
%     \item \textbf{Sentence transformers} - frozen pretrained encoders with a classification head, and a full reproduction of the published SetFit baseline.
%     \item \textbf{Ensembles} - soft voting over models with complementary per-project error profiles.
% \end{itemize}
\begin{table}[H]
  \centering
  \caption{Overview of evaluated models categorized by approach.}
  \label{tab:model-categories}
  \begin{tabular}{ll}
    \toprule
    \textbf{Category} & \textbf{Model} \\
    \midrule
    \textbf{Reference floors} & majority class \\
    & stratified random \\
    & hand-written keyword rules \\
    & naive Bayes implemented from scratch \\
    \addlinespace
    \textbf{Classical machine learning} & TF-IDF with naive Bayes \\
    & TF-IDF with logistic regression \\
    & TF-IDF with linear SVM \\
    & TF-IDF with random forest \\
    \addlinespace
    \textbf{Neural networks} & feed-forward network over TF-IDF features \\
    & 1-D convolutional network over learned embeddings \\
    \addlinespace
    \textbf{Sentence transformers} & frozen pretrained encoders with a classification head \\
    & full reproduction of the published SetFit baseline \\
    \addlinespace
    \textbf{Ensembles} & soft voting over models with complementary per-project error profiles \\
    \bottomrule
  \end{tabular}
\end{table}

\subsection{Structure}

Section 2 characterises the dataset. Section 3 describes our approach, including the persistence architecture and the preprocessing pipeline. Section 4 sets out the evaluation methodology. Section 5 presents results, Section 6 analyses where the best model fails, and Section 7 discusses threats to validity.